\subsection{Statements reproduced and statements predicted}
\label{sec:design-hypotheses}

Two kinds of statement are examined, and they differ in what an outcome could have done to them. The first kind, R1 to R4 in Table~\ref{tab:hypotheses}, are the statements of the design document of the product about the Annex 1 treatment: that the indexation charge is 55.9 percent of the delta NII of the short UF variant; that a single-currency treatment of the same flows gives a delta NII 6.8 to 8.8 times smaller and a delta EVE 10 to 27 percent smaller in the three variants whose net UF position is 30 percent of level 1 capital; and that no post-shock floor binds on the committed curves, the closest rate being 20.48 basis points from its floor. The case computed the documents on 16 October 2026 and the statements entered the design document on 21 October 2026, after the numbers were known, so no outcome of this paper could have refuted them. They are reproduced, with the code of this paper and not by reading the documents, and are not tested.

The second kind are five statements that the author wrote into the headers of the experiment scripts before the scripts that decide them were first run. The order cannot be shown by any history of the repository, because the manuscript and its scripts are committed together after the runs; the statements are the author's word. H1 is the bound of Corollary~\ref{cor:measures} applied to the data: for every variant, book and treatment, the sum of the legs is at most the pair rule, which is at most the positive parts. H2 predicts that, for the value measure, the pair rule has the largest absolute Shapley value in the three variants with a 30 percent net UF position. H3 states that the Annex 1 income measure exceeds the income measures of the single-currency treatment and of the Basel positive-part aggregation, and the Annex 1 value measure is not below the single-currency one, in every variant and book; it was written after the reproduction run had printed the committed-seed values, in which the CLP steepener variant already has the Annex 1 income measure below the positive-part one, so H3 is a statement that the data had already contradicted and is reported as such, not as a prediction. H4 states that there is a net UF position, within the range the generator can balance, at which the Annex 1 treatment flags the bank as prioritised in more than half of the books while the single-currency treatment flags none. H5 states that a binding UR floor never raises the Annex 1 measure. Each is refuted by a stated outcome: one violation of the bound for H1, a different largest switch in any of the three variants for H2, one exception for H3 and H5, and no such position for H4.

\begin{table}[htbp]
\centering
\footnotesize
\setlength{\tabcolsep}{3.5pt}
\caption{The statements tested or reproduced in this paper, with the deciding quantity and the outcome: R1 to R4 are the statements of the product's specification, entered after the numbers of the published case were known, and so reproduced and not tested; H1, H2, H4 and H5 were written by the author into the headers of the experiment scripts before the scripts that decide them ran (an order that the repository history cannot show); H3 was written after the reproduction run had printed a value that already contradicts one of its clauses, so it is a refuted statement and not a prediction.}
\label{tab:hypotheses}
\begin{tabular}{p{0.8cm}p{7.6cm}p{1.9cm}p{3.6cm}}
\toprule
Id & Statement and deciding quantity & Status & Outcome \\
\midrule
R1 & The indexation charge is 55.9\% of the delta NII of the short UF variant (reproduction). & reproduced & 55.9\% \\
R2 & The single-currency delta NII is 6.8 to 8.8 times smaller in the three 30\% variants (reproduction). & reproduced & 6.8 to 8.8 \\
R3 & The single-currency delta EVE is 10 to 27\% smaller in the same three variants (reproduction). & reproduced & 10 to 27\% \\
R4 & No floor binds on the committed curves; the closest post-shock rate is 20.48 bp from its floor (reproduction). & reproduced & 20.48 bp in every variant \\
H1 & The sum of the legs is at most the pair rule, which is at most the positive parts, for the two measures of every variant, book and treatment. & held & 0 violations in 2{,}880 comparisons \\
H2 & For the delta EVE, the pair rule has the largest absolute Shapley value in the three 30\% variants. & refuted & the own curve has the largest value in all three \\
H3 & The Annex 1 delta NII exceeds the single-currency and the positive-part delta NII, and the Annex 1 delta EVE is not below the single-currency one, in every variant and book. & refuted in one clause (not a prediction) & the CLP steepener variant: Annex 1 below positive parts in 30 of 30 books \\
H4 & A net UF position exists at which Annex 1 flags the bank in most books and the single-currency treatment flags none. & held at one position only & 110\% of Tier 1 only (17 of 20 books against 0), the delta EVE test only \\
H5 & A binding UR floor never raises the Annex 1 measure. & refuted & 208 of 1{,}920 comparisons of the two measures raise it, all on the delta NII \\
\bottomrule
\end{tabular}
\end{table}


\subsection{Data and provenance}
\label{sec:design-data}

The curves are real. The Banco Central de Chile publishes in table F022 of its statistical database the daily secondary-market yields of Central Bank and Treasury bonds, in pesos at 1, 2, 5 and 10 years and in UF at 1, 2, 5, 10, 20 and 30 years, together with the UF value and the monetary policy rate \citep{BCCh_BDE}; the terms of use allow the information to be reproduced and adapted when the central bank is named as the owner, and every table derived from it here is an adaptation by the author. The raw files of September 2026 hold \nTenorsClp{} peso points, of which three are yields (2, 5 and 10 years; the one-year series has no observation) and the fourth is the policy rate at one day, which the engine's four-parameter fit requires, and \nTenorsUr{} UF points (2, 5, 10 and 30 years; the 1-year and 20-year series have none). The valuation date is \valDate{}, the last day of the month on which every chosen series has a value, and the UF value of that day is \ufValue{} pesos. The fitted zero curves have root-mean-square errors of \rmseClp{} and \rmseUr{} basis points; the peso curve runs from \shortClp{} to \longClp{} percent over the 19 midpoints and the UF curve from \shortUr{} to \longUr{} percent. Table~\ref{tab:data} lists the inputs.

\begin{table}[htbp]
\centering
\footnotesize
\setlength{\tabcolsep}{3.5pt}
\caption{The inputs of the experiments and where each comes from: the daily secondary-market yields of the F022 table of the Banco Central de Chile on the valuation date 2026-09-03, the UF value of that day, the synthetic bank built by the case's generator and the rules of Annex 1 that the engine holds.}
\label{tab:data}
\begin{tabular}{p{2.4cm}p{6.6cm}p{5.4cm}}
\toprule
Input & Source and identifier & Use in the experiments \\
\midrule
Peso curve & Banco Central de Chile, table F022 (secondary-market yields), series BCLP at 2, 5, 10 years, with the monetary policy rate (series TPM) as the overnight point & zero curve of the CLP leg, 19 midpoints \\
UF curve & Banco Central de Chile, table F022, series BUF at 2, 5, 10, 30 years & zero curve of the UR leg, 19 midpoints \\
UF value & BDE series F073.UFF.PRE.Z.D, value of 2026-09-03 & conversion of UF flows to pesos \\
Balance sheet & synthetic Banco Andino, calibrated to the CMF aggregates of August 2026 & 30 seeds of the repricing profiles per variant \\
Shock sizes, floors, offset, charge & RAN 21-13 Annex 1, Hojas 9 to 12 (held by the engine) & all treatments \\
\bottomrule
\end{tabular}
\end{table}


The balance sheet is synthetic. Banco Andino is generated by the product's own generator: its totals, the loan mix, the securities, deposits and equity are calibrated to the aggregates that the supervisor publishes for the banking system for August 2026, level 1 capital is \tierOne{} million pesos, the peso leg has \clpAssets{} million pesos of assets and \clpLiabs{} of liabilities in \nItemsClp{} items, and the UF leg \urAssetsUf{} million UF of assets and \urLiabsUf{} of liabilities in \nItemsUr{} items. The UF share of each portfolio, the repricing profile of each item, the deposit split and the behavioural parameters are synthetic and stated, and a seed perturbs the repricing profile of every item by a log-normal factor, renormalised, so that two seeds give two banks of the same size and totals. The UF borrowings are solved so that the net UF position of a variant is its stated share of level 1 capital. The variants are the six of the case: long UF and short UF (net position of $\pm30$ percent), a long UF book under a surprise of 2.5 percentage points in inflation, a CLP steepener (5 percent, with 35 percent of the fixed peso loans and securities moved to the long end), offset on and off (30 percent, with the peso bonds at twice the maturity) and a floor-exposed book (5 percent, with 70 percent of the UF commercial loans floating).

\subsection{The methods compared, and how they are run}
\label{sec:design-methods}

Every method compared is run by the product on the same flows; no number is taken from a publication. The Annex 1 treatment is the case with all five switches on. The single-currency treatment is the case with all off. The Basel positive-part aggregation is run with the engine's \texttt{eve\_risk\_measure} on the matrix of leg deltas of the Annex 1 curves, sizes and floors, and without a charge, because the Basel standard has no indexation charge. The repricing gap and the duration gap are run by the engine on the whole book as one currency. The published yield-curve models of Section~\ref{sec:rw-curves} enter only through the engine's bootstrap and Nelson-Siegel fit; the Vasicek, Hull-White and Heath-Jarrow-Morton families are not run, because they generate paths for an earnings-at-risk measure and the standardised measures of the supervisor need only the curve at the 19 midpoints. The engine's earnings-at-risk simulation is not run either, for the same reason, and the case's committed simulation of the Hull-White one-factor paths is not used.

\subsection{Metrics, statistics and protocol}
\label{sec:design-protocol}

The metrics are the two measures of Equations~\eqref{eq:measure-eve} and~\eqref{eq:measure-nii} and their analogues under a treatment, Equation~\eqref{eq:measure-tau}, in millions of pesos and as a share of level 1 capital. The relative shortfall of the single-currency treatment is
\begin{equation}
\sigma_{\mathrm{EVE}}=100\Bigl(1-\frac{M_{\mathrm{EVE}}(\mathbf 0)}{M_{\mathrm{EVE}}(\mathbf 1)}\Bigr)\ \text{percent},\qquad
\rho_{\mathrm{NII}}=\frac{M_{\mathrm{NII}}(\mathbf 1)}{M_{\mathrm{NII}}(\mathbf 0)},
\label{eq:shortfall}
\end{equation}
and the share of the charge in the income measure is $\kappa|\PNUR|/M_{\mathrm{NII}}(\mathbf 1)$. The main comparison runs \mainSeeds{} books per variant, seeds 1 to 30 of the repricing profiles, on the curves of \valDate{}; the committed seed 42 is run in addition and reproduces the committed documents. Every number of the comparison is the mean over the books with the half-width of its 95 percent Student interval, the standard deviation over the books divided by the square root of 30 times the quantile of Student's distribution with 29 degrees of freedom. The curves and the UF value do not vary across books, so the interval measures the variation of the synthetic balance sheet and not of the market. The sweep over the net UF position uses \sweepSeeds{} books (seeds 101 to 120) at each position from $-150$ to $+150$ percent of level 1 capital in steps of 10 points; the generator refuses the positions that it cannot balance, and they are listed. The floor experiment uses \floorsSeeds{} books (seeds 201 to 220) at each translation of the UF curve of 0, $-25$, $-50$, $-100$, $-150$, $-200$, $-250$ and $-300$ basis points, a translation that leaves the peso curve alone and is a what-if on the level of the real curve and not a historical curve: the raw files hold a single month, and the central bank's interface is not queried. A book is flagged as prioritised when the engine's \texttt{prioritised\_bank} reports a breach of any of the three tests of the text. The decomposition of Equation~\eqref{eq:shapley} is computed exactly from the 32 treatments, and the check that its five values add up to the difference of the measures is asserted for every book.

\subsection{Ablations planned}
\label{sec:design-ablation}

The ablation of Section~\ref{sec:ablation} is the five-switch factorial itself. It reports, for the long UF variant, the change of each measure when a switch is turned off in Annex 1, the change when it alone is turned on in the single-currency treatment, and the Shapley value; and, for the three variants with a 30 percent position, the Shapley values with their intervals. The difference between the first two quantities for one switch, and its distance from the Shapley value, measure how much that switch interacts with the others. The ablation does not search over other components: the five switches are the five inputs in which the two treatments differ, and the other rules of the text (the grid, the six scenarios, the behavioural slotting of the deposits, the prepayment and redemption factors) are common to both treatments and held fixed. The bound of Proposition~\ref{prop:pair} is checked on the eight combinations of curve, sizes and floor with the pair switch on and the charge off, so that it is tested on treatments that the decomposition visits and on others it does not report.
